Building upon The fundamental principles of neural symbols are embodied in the knowledge graph-centric construction approach of in M$^3$-Med \cite{m3med}, we introduce \PipelineName, a novel framework that formalizes benchmark synthesis as a knowledge-driven engineering process. 
% 基于 M$^3$-Med \cite{m3med} 中以知识图谱为中心的构建方法所体现的神经符号的基本原理，我们引入了 \PipelineName，这是一个将基准合成形式化为知识驱动的工程过程的新框架。
Fundamentally diverging from simulation engines that lack medical fidelity and stochastic LLM approaches prone to hallucination, we reframe the task of dataset construction as a deterministic graph traversal problem over structured video representations. 
% 与缺乏生物学保真度的仿真引擎和容易产生幻觉的随机LLM方法截然不同，我们将数据集构建任务重新定义为在结构化视频表示上进行确定性图遍历的问题。
By explicitly anchoring detected visual primitives within a fine-grained Medical Knowledge Graph, our system enforces a verifiable, deterministic Chain-of-Thought (CoT) for every generated query. 
% 通过将检测到的视觉基元显式地锚定在细粒度的医学知识图谱中，我们的系统为每个生成的查询强制执行可验证的确定性思维链（CoT）。
This architectural design enables the systematic encoding of complex, multi-hop reasoning logic—spanning causal inference and temporal sequencing—into benchmark queries, while ensuring full logical provenance and interpretability.
% 这种架构设计能够将复杂的、多跳推理逻辑（涵盖因果推理和时间序列）系统地编码到基准查询中，同时确保完整的逻辑溯源性和可解释性。

Our architectural workflow decomposes the unstructured video processing into three tightly coupled abstraction layers:
% 我们的架构工作流程将非结构化视频处理分解为三个紧密耦合的抽象层：
\begin{itemize}
    \item \textbf{The Visual Extraction Layer} (Pixel-Level): This module ingests raw video data to extract distinct visual primitives. By analyzing spatiotemporal proximity and feature vector similarity, it associates discrete detections into continuous 3D object trajectories, formally termed as \textit{Spatiotemporal Tubelets}.
    \item \textbf{The Graph Construction Layer} (Semantic-Level): These tubelets are instantiated as entities within a dynamic Knowledge Graph. We quantify pairwise spatiotemporal and semantic affinities to establish logic-weighted edges, effectively transforming the video into a structured property graph.
    \item \textbf{The Query Synthesis Layer} (Logic-Level): By executing deterministic traversal algorithms (e.g., DFS) over the graph, we extract valid reasoning paths. These paths are then expanded via generative models into natural language questions, with the traversal endpoints serving as ground-truth video segments.
\end{itemize}
% 1. 视觉提取层（像素级）：该模块接收原始视频数据以提取不同的视觉基元。通过分析时空邻近性和特征向量相似性，它将离散的检测结果关联到连续的 3D 对象轨迹，正式称为“时空管状结构”。
% 2. 图构建层（语义级）：这些管状结构被实例化为动态知识图谱中的实体。我们量化成对的时空和语义亲和性，以建立逻辑加权边，从而有效地将视频转换为结构化的属性图。
% 3. 查询合成层（逻辑级）：通过在图上执行确定性遍历算法（例如，深度优先搜索），我们提取有效的推理路径。然后通过生成模型将这些路径扩展为自然语言问题，遍历终点作为真实视频片段。



